%======================================================================
\section{From the QCD Lagrangian to the light-cone QCD Hamiltonian}
\label{app:hamiltonian}
%======================================================================

The QCD Lagrangian with $\Nf$ massless flavours is
\begin{equation}
\mathcal L_{\rm QCD}=-\frac14F^{a}_{\mu\nu}F^{a\mu\nu}+\bar\psi\,i\gamma^{\mu}D_{\mu}\psi,
\qquad D_{\mu}=\partial_{\mu}-igA^{a}_{\mu}t^{a}.
\label{eq:LQCD}
\end{equation}
The Yang--Mills part is treated as in appendix~A of Paper~I; the only change is that the colour
current entering Gauss's law is now $J^{a\nu}=\bar\psi\gamma^{\nu}t^{a}\psi$ plus the gluon
current. In light-cone gauge $A^{+}=0$ and with $\psi_{\pm}=\Lambda_{\pm}\psi$, the Dirac equation
splits into a dynamical equation for $\psi_{+}$ and a constraint for $\psi_{-}$,
\begin{equation}
i\partial^{+}\psi_{-}=\frac12\gamma^{0}\big(i\gamma^{i}D_{i}\big)\psi_{+}
\quad\Rightarrow\quad
\psi_{-}=\frac{1}{2i\partial^{+}}\,\gamma^{0}\gamma^{i}\big(i\partial_{i}+gA_{i}\big)\psi_{+},
\label{eq:psiminus}
\end{equation}
while the $+$ component of the Yang--Mills equation determines $A^{-}$,
\begin{equation}
\partial^{+}A^{a-}=-\frac{1}{\partial^{+}}\Big(D^{ab}_{i}\partial^{+}A^{b}_{i}-gJ^{a+}\Big),
\qquad J^{a+}=\bar\psi\gamma^{+}t^{a}\psi=\sqrt2\,\psi_{+}^{\dg}t^{a}\psi_{+}.
\label{eq:Aminus}
\end{equation}
Eliminating $\psi_{-}$ and $A^{-}$ through \eqref{eq:psiminus} and \eqref{eq:Aminus} from the
canonical Hamiltonian $P^{-}=\int dx^{-}d^{2}x\,[\Pi\,\partial^{-}\phi-\mathcal L]$, one obtains
\cite{Brodsky:1997de,Zhang:1993dd,Kogut:1969xa} the light-cone Hamiltonian
\eqref{eq:HQCD}--\eqref{eq:Hq}, in which the fields are free fields at $x^{+}=0$. The term
$\bar\psi\gamma^{i}A_{i}\,(\gamma^{+}/2i\partial^{+})\,\gamma^{j}A_{j}\psi$ comes from the
$gA_{i}$ part of $\psi_{-}$ and describes the instantaneous fermion propagation; the current--current
terms come from the $gJ^{a+}$ part of $A^{-}$ and combine with the gluonic current of Paper~I into
the single operator \eqref{eq:HinstA}. In the two-component representation
$\psi_{+}\leftrightarrow\xi$ of Refs.~\cite{Zhang:1993dd,Brodsky:1997de}, the quark--gluon vertex
becomes the bilinear \eqref{eq:Phi} between the two-component spinors, the factors $2k^{i}/k^{+}$,
$(\bm\sigma\!\cdot\!\bm P_{L})\sigma^{i}/P^{+}_{L}$ and $\sigma^{i}(\bm\sigma\!\cdot\!\bm P_{R})/P^{+}_{R}$
arising from $\partial_{i}A_{i}/\partial^{+}$ and from the two $\bm\sigma\!\cdot\!\bm\partial/\partial^{+}$
operators of \eqref{eq:psiminus} acting on the outgoing and incoming fermion lines.

%======================================================================
\section{State definitions and normalization}
\label{app:states}
%======================================================================

As in Paper~I, $\vac$ denotes the vacuum of the free soft Hamiltonian, $\vac\equiv\vac_{\rm soft}\otimes|v\rangle$,
with $E_{0}=0$. The single particle states are
\begin{equation}
\ket{\gl{a}{i}{k}}=\frac{a^{\dg a}_{i}(k)}{(2\pi)^{3/2}}\vac,\qquad
\ket{\qk{\beta}{\lambda}{l}}=\frac{b^{\dg\beta}_{\lambda}(l)}{(2\pi)^{3/2}}\vac,\qquad
\ket{\aqk{\alpha}{\lambda}{m}}=\frac{d^{\dg\alpha}_{\lambda}(m)}{(2\pi)^{3/2}}\vac,
\label{eq:states1}
\end{equation}
normalized as
$\langle\qk{\beta'}{\lambda'}{l'}|\qk{\beta}{\lambda}{l}\rangle=\delta^{\beta\beta'}\delta_{\lambda\lambda'}\dthree(l-l')$,
with energies $\eps_{k}$, $\eps_{l}$, $\eps_{m}$. Multiparticle states are defined by the product of
creation operators, each with a factor $(2\pi)^{-3/2}$, in the order
\begin{equation}
\ket{\qk{\beta}{\lambda_{1}}{l}\aqk{\alpha}{\lambda_{2}}{m}\,g\cdots}=
\frac{b^{\dg\beta}_{\lambda_{1}}(l)\,d^{\dg\alpha}_{\lambda_{2}}(m)\,a^{\dg}\cdots}{(2\pi)^{3n/2}}\vac ,
\label{eq:statesn}
\end{equation}
quark operators to the left of antiquark operators; gluon operators commute with both. The energy of
a multiparticle state is the sum of the single particle energies. With the ordering
\eqref{eq:statesn} and the ordering $d\,b$ of the pair annihilation operators in \eqref{eq:Om2qq},
\begin{equation}
d^{\alpha}_{\lambda_{2}}(m)\,b^{\beta}_{\lambda_{1}}(l)\;b^{\dg\beta'}_{\lambda_{1}'}(l')\,d^{\dg\alpha'}_{\lambda_{2}'}(m')\vac
=(2\pi)^{6}\delta^{\beta\beta'}\delta^{\alpha\alpha'}\delta_{\lambda_{1}\lambda_{1}'}\delta_{\lambda_{2}\lambda_{2}'}
\dthree(l-l')\dthree(m-m')\vac,
\end{equation}
with a positive sign, so that all the matrix elements of $\OM$ between states with one soft pair are
read off without additional signs. In states with two identical fermions the antisymmetry of the
creation operators is accounted for by the explicit antisymmetrization of the corresponding
coefficients, Eqs.~\eqref{eq:UK2} and \eqref{eq:K1}.

The matching rule \eqref{eq:rule} follows from these definitions. A normal ordered term with
coefficient $K$, $M$ creation and $n$ annihilation operators, acting on an $n$-particle state
contracts each annihilation operator with an incoming particle, which produces a factor
$(2\pi)^{3}\delta^{(3)}$ against the measure $1/(2\pi)^{3}$ and the state normalization
$(2\pi)^{-3/2}$, i.e.\ a net $(2\pi)^{-3/2}$ per incoming particle; each of the $M$ creation operators
comes with the measure $1/(2\pi)^{3}$ and builds a normalized outgoing particle with
$(2\pi)^{3/2}$, i.e.\ again $(2\pi)^{-3/2}$. Hence $\OM\ket{n}\supset(2\pi)^{-3(n+M)/2}\int d^{3}(M)\,K\ket{M}$.

The notation for the wave functions follows Paper~I: the subscript lists the outgoing particles and
the powers of $\rho$, the argument the incoming particles, and all outgoing and intermediate momenta
are integrated over.

%======================================================================
\section{The CGC Hamiltonian with quarks: eikonal approximation}
\label{app:eikonal}
%======================================================================

The soft and valence modes of the quark field are introduced as in \eqref{eq:split},
\begin{equation}
\xi=\Xi+\xi,\qquad
\Xi^{\beta}(x)=2^{-1/4}\!\int_{\vee}^{\infty}\!\frac{dk^{+}}{2\pi}\!\int\!\frac{d^{2}k}{(2\pi)^{2}}
\sum_{\lambda}\Big[\bar b^{\beta}_{\lambda}(k)\chi_{\lambda}e^{-ik\cdot x}+\bar d^{\dg\beta}_{\lambda}(k)\chi_{\lambda}e^{ik\cdot x}\Big],
\end{equation}
with the soft field $\xi$ restricted to $\Lambda<k^{+}<\vee$, and similarly for the gluon field
(appendix~C of Paper~I). Inserting the decomposition into \eqref{eq:Hq} generates soft--soft,
soft--valence and valence--valence terms. The eikonal approximation neglects all contributions
suppressed by powers of $k^{+}_{\rm soft}/k^{+}_{\rm valence}$ and all terms in which inverse longitudinal
derivatives act on valence fields. Concretely:
\begin{itemize}
\item the vertex $\bar\psi\gamma^{i}t^{a}\psi A^{a}_{i}$ with a valence fermion and a soft gluon changes
the valence momentum by a soft amount; its transverse structure \eqref{eq:Phi} contains
$1/P^{+}_{L,R}$ of the valence fermion and the only unsuppressed part is the $2k^{i}/k^{+}$ term,
which is the coupling of the soft gluon to the plus current of the valence quark. Together with the
corresponding term of the gluonic sector it builds $H_{g}$ with the total charge density
\eqref{eq:rho};
\item vertices with a soft fermion and a valence gluon, or with one soft and one valence fermion, cannot
conserve the plus momentum within the soft and valence windows and are dropped;
\item the instantaneous fermion exchange with valence fields carries $1/P^{+}$ of a valence line and is
suppressed;
\item in the instantaneous gluon exchange \eqref{eq:HinstA}, the valence current reduces to its colour
charge, $J^{a+}_{\rm val}\to\rho^{a}$, and the cross terms with the soft currents give the
$\rho\cdot J_{g}$ term of Paper~I and the $\rho\cdot J_{q}$ term of \eqref{eq:Hqinst}.
\end{itemize}
The valence quark charge density \eqref{eq:rhoq} is the Fourier transform of
$J^{a+}_{\rm val,q}=\sqrt2\,\Xi^{\dg}t^{a}\Xi$, normal ordered, restricted to its colour rotating
($b^{\dg}b$ and $d^{\dg}d$) part; the terms that create or annihilate valence pairs change the valence
plus momentum by a valence amount and do not contribute. Using \eqref{eq:bb} for the valence modes
one verifies \eqref{eq:rhoalg} for $\rho_{q}$; since $\rho_{g}$ and $\rho_{q}$ commute, the algebra holds for
their sum.

%======================================================================
\section{Matrix elements}
\label{app:me}
%======================================================================

We list the matrix elements used in the calculation, between the states of
appendix~\ref{app:states}; colour indices in the fundamental representation are $\alpha,\beta,\dots$,
spin labels $\lambda$, and $\Ph$ is defined in \eqref{eq:Phi}.

\noindent$\bullet$ Gluon emission and absorption by the valence charge (Paper~I), with spectators $X$:
{\small\begin{align}
\bra{\gl{a}{i}{k}}H_{g}\vac&=\frac{g\,k^{i}\rho^{a}(-k)}{4\pi^{3/2}|k^{+}|^{3/2}},
\label{eq:MEg}\\
\bra{X\,\gl{b}{j}{p}}H_{g}\ket{X}=\frac{g\,p^{j}\rho^{b}(-p)}{4\pi^{3/2}|p^{+}|^{3/2}},&\qquad
\bra{X}H_{g}\ket{X\,\gl{b}{j}{p}}=\frac{g\,p^{j}\rho^{b}(p)}{4\pi^{3/2}|p^{+}|^{3/2}} .
\label{eq:MEg2}
\end{align}}
$\bullet$ Three-gluon vertex (Paper~I):
{\small\begin{equation}
\bra{\gl{c}{l}{q}\gl{b}{j}{p}}H_{ggg}\ket{\gl{a}{i}{k}}=
\frac{igf^{abc}\dthree(k-p-q)}{8\pi^{3/2}\sqrt{k^{+}p^{+}q^{+}}}\Vg^{ijl}(k;p,q).
\label{eq:MEggg}
\end{equation}}
$\bullet$ Quark--gluon vertices:
{\small\begin{align}
\bra{\qk{\beta}{\lambda_{1}}{l}\aqk{\alpha}{\lambda_{2}}{m}}H_{gq\bar q}\ket{\gl{a}{i}{k}}&=
\tfrac{g\,t^{a}_{\beta\alpha}}{8\pi^{3/2}\sqrt{k^{+}}}\,\dthree(k-l-m)\,\Ph^{i}_{\lambda_{1}\lambda_{2}}(k;l,m),
\label{eq:MEqqg}\\
\bra{\qk{\beta}{\lambda}{l}\gl{b}{j}{p}}H_{gq\bar q}\ket{\qk{\beta'}{\lambda'}{l'}}&=
\tfrac{g\,t^{b}_{\beta\beta'}}{8\pi^{3/2}\sqrt{p^{+}}}\,\dthree(l'-l-p)\,\Ph^{j}_{\lambda\lambda'}(p;l,l'),
\label{eq:MEqg}\\
\bra{\aqk{\alpha}{\lambda}{m}\gl{b}{j}{p}}H_{gq\bar q}\ket{\aqk{\alpha'}{\lambda'}{m'}}&=
-\tfrac{g\,t^{b}_{\alpha'\alpha}}{8\pi^{3/2}\sqrt{p^{+}}}\,\dthree(m'-m-p)\,\Ph^{j}_{\lambda'\lambda}(p;m',m).
\label{eq:MEaqg}
\end{align}}
The absorption vertices are the complex conjugates, with
$\big[\Ph^{i}_{\lambda\lambda'}(k;P_{L},P_{R})\big]^{*}=\Ph^{i}_{\lambda'\lambda}(k;P_{R},P_{L})$.

\noindent$\bullet$ Instantaneous interaction of soft quarks with the valence charge ($\rho\cdot J_{q}$):
{\small\begin{align}
\bra{\qk{\beta}{\lambda_{1}}{l}\aqk{\alpha}{\lambda_{2}}{m}}H_{q-{\rm inst}}\vac&=
\frac{g^{2}\,t^{a}_{\beta\alpha}\,\rho^{a}(-l-m)\,\delta_{\lambda_{1}\lambda_{2}}}{8\pi^{3}(l^{+}+m^{+})^{2}},
\label{eq:MEinst}\\
\bra{\qk{\beta}{\lambda}{l}}H_{q-{\rm inst}}\ket{\qk{\beta'}{\lambda'}{l'}}&=
\frac{g^{2}\,t^{a}_{\beta\beta'}\,\rho^{a}(l'-l)\,\delta_{\lambda\lambda'}}{8\pi^{3}(l'^{+}-l^{+})^{2}},
\label{eq:MEinstq}\\
\bra{\aqk{\alpha}{\lambda}{m}}H_{q-{\rm inst}}\ket{\aqk{\alpha'}{\lambda'}{m'}}&=
-\frac{g^{2}\,t^{a}_{\alpha'\alpha}\,\rho^{a}(m'-m)\,\delta_{\lambda\lambda'}}{8\pi^{3}(m'^{+}-m^{+})^{2}}.
\label{eq:MEinstqb}
\end{align}}
$\bullet$ Instantaneous quark exchange ($H_{ggq\bar q}$); the instantaneous fermion line carries
$1/(2P^{+})$, $P^{+}$ being the plus momentum flowing along the fermion number arrow:
{\small\begin{align}
\bra{\qk{\beta}{\lambda_{1}}{l}\aqk{\alpha}{\lambda_{2}}{m}\gl{d}{n}{r}}H_{ggq\bar q}\ket{\gl{a}{i}{k}}&=
\frac{g^{2}\,\dthree(k-l-m-r)}{16\pi^{3}\sqrt{k^{+}r^{+}}}\,
\chi^{\dg}_{\lambda_{1}}\!\Big[\frac{(t^{d}t^{a})_{\beta\alpha}\sigma^{n}\sigma^{i}}{l^{+}+r^{+}}
-\frac{(t^{a}t^{d})_{\beta\alpha}\sigma^{i}\sigma^{n}}{m^{+}+r^{+}}\Big]\chi_{\lambda_{2}},
\label{eq:MEggqq1}\\
\bra{\qk{\beta}{\lambda_{1}}{l}\aqk{\alpha}{\lambda_{2}}{m}}H_{ggq\bar q}\ket{\gl{a}{i}{k}\gl{b}{j}{p}}&=
\frac{g^{2}\,\dthree(k+p-l-m)}{16\pi^{3}\sqrt{k^{+}p^{+}}}\,
\chi^{\dg}_{\lambda_{1}}\!\Big[\frac{(t^{a}t^{b})_{\beta\alpha}\sigma^{i}\sigma^{j}}{l^{+}-k^{+}}
+\frac{(t^{b}t^{a})_{\beta\alpha}\sigma^{j}\sigma^{i}}{l^{+}-p^{+}}\Big]\chi_{\lambda_{2}},
\label{eq:MEggqq2}\\
\bra{\qk{\beta}{\lambda}{l}\gl{c}{k}{q}\gl{d}{n}{r}}H_{ggq\bar q}\ket{\qk{\beta'}{\lambda'}{l'}}&=
\frac{g^{2}\,\dthree(l'-l-q-r)}{16\pi^{3}\sqrt{q^{+}r^{+}}}\,
\chi^{\dg}_{\lambda}\!\Big[\frac{(t^{c}t^{d})_{\beta\beta'}\sigma^{k}\sigma^{n}}{l^{+}+q^{+}}
+\frac{(t^{d}t^{c})_{\beta\beta'}\sigma^{n}\sigma^{k}}{l^{+}+r^{+}}\Big]\chi_{\lambda'},
\label{eq:MEq2g}\\
\bra{\qk{\beta}{\lambda}{l}\gl{d}{n}{r}}H_{ggq\bar q}\ket{\qk{\beta'}{\lambda'}{l'}\gl{c}{k}{q}}&=
\frac{g^{2}\,\dthree(l'+q-l-r)}{16\pi^{3}\sqrt{q^{+}r^{+}}}\,
\chi^{\dg}_{\lambda}\!\Big[\frac{(t^{d}t^{c})_{\beta\beta'}\sigma^{n}\sigma^{k}}{l'^{+}+q^{+}}
+\frac{(t^{c}t^{d})_{\beta\beta'}\sigma^{k}\sigma^{n}}{l'^{+}-r^{+}}\Big]\chi_{\lambda'}.
\label{eq:MEqginst}
\end{align}}
$\bullet$ Instantaneous gluon exchange between soft currents ($J_{g}\cdot J_{q}$ and $J_{q}\cdot J_{q}$):
{\small\begin{align}
\bra{\qk{\beta}{\lambda_{1}}{l}\aqk{\alpha}{\lambda_{2}}{m}\gl{d}{n}{r}}H_{q-{\rm inst}}\ket{\gl{a}{i}{k}}&=
\frac{ig^{2}f^{bad}\,t^{b}_{\beta\alpha}\,\delta_{\lambda_{1}\lambda_{2}}\delta_{in}\,(k^{+}+r^{+})\,\dthree(k-l-m-r)}
{2(2\pi)^{3}\sqrt{k^{+}r^{+}}\,(l^{+}+m^{+})^{2}},
\label{eq:MEJJ1}\\
\bra{\qk{\beta}{\lambda_{1}}{l}\aqk{\alpha}{\lambda_{2}}{m}}H_{q-{\rm inst}}\ket{\gl{a}{i}{k}\gl{b}{j}{p}}&=
-\frac{ig^{2}f^{cab}\,t^{c}_{\beta\alpha}\,\delta_{\lambda_{1}\lambda_{2}}\delta_{ij}\,(p^{+}-k^{+})\,\dthree(k+p-l-m)}
{2(2\pi)^{3}\sqrt{k^{+}p^{+}}\,(k^{+}+p^{+})^{2}},
\label{eq:MEJJ2}\\
\bra{\qk{\beta}{\lambda}{l}\gl{c}{k}{q}\gl{d}{n}{r}}H_{q-{\rm inst}}\ket{\qk{\beta'}{\lambda'}{l'}}&=
\frac{ig^{2}f^{acd}\,t^{a}_{\beta\beta'}\,\delta_{\lambda\lambda'}\delta_{kn}\,(r^{+}-q^{+})\,\dthree(l'-l-q-r)}
{2(2\pi)^{3}\sqrt{q^{+}r^{+}}\,(q^{+}+r^{+})^{2}},
\label{eq:MEJJq2g}\\
\bra{\qk{\beta}{\lambda}{l}\gl{d}{n}{r}}H_{q-{\rm inst}}\ket{\qk{\beta'}{\lambda'}{l'}\gl{c}{k}{q}}&=
\frac{ig^{2}f^{acd}\,t^{a}_{\beta\beta'}\,\delta_{\lambda\lambda'}\delta_{kn}\,(q^{+}+r^{+})\,\dthree(l'+q-l-r)}
{2(2\pi)^{3}\sqrt{q^{+}r^{+}}\,(q^{+}-r^{+})^{2}},
\label{eq:MEJJqg}\\
\bra{\qk{\beta}{\lambda_{1}}{l}\aqk{\alpha}{\lambda_{2}}{m}}H_{q-{\rm inst}}\ket{\qk{\beta'}{\lambda_{1}'}{l'}\aqk{\alpha'}{\lambda_{2}'}{m'}}&
\nonumber\\
=\frac{g^{2}\,\dthree(l+m-l'-m')}{(2\pi)^{3}}\Big[&
\frac{t^{a}_{\beta\alpha}t^{a}_{\alpha'\beta'}\delta_{\lambda_{1}\lambda_{2}}\delta_{\lambda_{1}'\lambda_{2}'}}{(l^{+}+m^{+})^{2}}
-\frac{t^{a}_{\beta\beta'}t^{a}_{\alpha'\alpha}\delta_{\lambda_{1}\lambda_{1}'}\delta_{\lambda_{2}\lambda_{2}'}}{(l'^{+}-l^{+})^{2}}\Big],
\label{eq:MEJJqq}\\
\bra{\qk{\beta}{\lambda}{l}\qk{\beta_{2}}{\lambda_{2}}{l_{2}}\aqk{\alpha}{\lambda_{m}}{m}}H_{q-{\rm inst}}\ket{\qk{\beta'}{\lambda'}{l'}}&
\nonumber\\
=\frac{g^{2}\,\dthree(l'-l-l_{2}-m)}{(2\pi)^{3}}\Big[&
\frac{t^{a}_{\beta\beta'}t^{a}_{\beta_{2}\alpha}\delta_{\lambda\lambda'}\delta_{\lambda_{2}\lambda_{m}}}{(l_{2}^{+}+m^{+})^{2}}
-\big((\beta,\lambda,l)\leftrightarrow(\beta_{2},\lambda_{2},l_{2})\big)\Big].
\label{eq:MEJJqqq}
\end{align}}
The remaining four-point matrix elements of the pure gauge sector are given in appendix~D of
Paper~I. Each instantaneous gluon line carries $1/(P^{+})^{2}$ with $P^{+}$ the plus momentum it
transfers; the relative minus signs reflect the antiquark charge and Fermi statistics.

%======================================================================
\section{Perturbation theory with non-commutative matrix elements}
\label{app:pt}
%======================================================================

The derivation of the wave functions through order $g^{2}$ follows the Rayleigh--Schr\"odinger
expansion of appendix~E of Paper~I \cite{Landau:1991,Shankar:2012}, keeping the order of the operator
valued matrix elements, the later interaction standing on the left. The presence of fermions only
introduces the signs dictated by the anticommutators \eqref{eq:bb} and the state ordering
\eqref{eq:statesn}. Here we prove the two identities used throughout.

\paragraph{Reduction formula.} With the $i\epsilon$ prescription of section~\ref{sec:lcpt},
$H_{mi}=(E_{i}-E_{m}+i\epsilon)\,\mu_{1}(m\leftarrow i)$ and $H_{fm}=(E_{m}-E_{f}+i\epsilon)\,\mu_{1}(f\leftarrow m)$.
Substituting into the second order amplitude \eqref{eq:mu12} gives \eqref{eq:reduction} identically.
For two subprocesses $X$, $Y$ acting on disjoint particles there are two intermediate states, with
$E_{m}-E_{f}=\Delta_{Y}$ ($X$ first) and $\Delta_{X}$ ($Y$ first), and $E_{i}-E_{f}=\Delta_{X}+\Delta_{Y}$; this
gives \eqref{eq:factor}. If neither or only one of $X$, $Y$ contains $\rho$, the two products coincide and
the overall denominator cancels. In operator language, the same statement follows from
$\OM=e^{iG}$: writing $\OM^{(2)}=iG^{(2)}+\tfrac12(\OM^{(1)})^{2}$ with
$iG^{(2)}_{fi}=\mathcal K_{fi}/(E_{i}-E_{f})$ (section~\ref{sec:diag}) and
$\mathcal K=H^{(2)}+\tfrac12[H^{(1)},\OM^{(1)}]$, one finds for a single intermediate state, with
$a=E_{i}-E_{m}$ and $b=E_{m}-E_{f}$,
\begin{equation}
\OM^{(2)}_{fi}=H_{fm}H_{mi}\Big[\frac{1}{2ab}+\frac{1}{2(a+b)}\Big(\frac1a-\frac1b\Big)\Big]
+\frac{H^{(2)}_{fi}}{a+b}
=\frac{H_{fm}H_{mi}}{a\,(a+b)}+\frac{H^{(2)}_{fi}}{a+b},
\end{equation}
which is the LCPT amplitude $\mu_{2}$. Thus the normal ordered coefficients of $\OM$ at second order are
the LCPT amplitudes, and the diagonalization of section~\ref{sec:diag} and the matching of
section~\ref{sec:coeff} are the same calculation.

\paragraph{Unitarity identity.} For a given intermediate state $m$ let $a=E_{i}-E_{m}+i\epsilon$ and
$b=E_{f}-E_{m}-i\epsilon$, so that $E_{i}-E_{f}+2i\epsilon=a-b$. Then
\begin{align}
\mu_{2}(f\leftarrow i)+\big[\mu_{2}(i\leftarrow f)\big]^{\dg}
&=H_{fm}H_{mi}\Big[\frac{1}{(a-b)\,a}-\frac{1}{(a-b)\,b}\Big]
\nonumber\\
&=-\frac{H_{fm}H_{mi}}{a\,b}
=-\big[\mu_{1}(m\leftarrow f)\big]^{\dg}\mu_{1}(m\leftarrow i),
\end{align}
which is \eqref{eq:unitid}; the instantaneous terms are anti-Hermitian,
$H^{(2)}_{fi}/(E_{i}-E_{f}+i\epsilon)+\big[H^{(2)}_{if}/(E_{f}-E_{i}+i\epsilon)\big]^{\dg}=0$ up to terms
proportional to $\epsilon$, and drop out. The ordering $H_{fm}H_{mi}$ is preserved, so the identity
holds for the non-commuting coefficients. Every unitarity relation of section~\ref{sec:omega}
is a sum of such identities over the intermediate states.
